\documentclass[aspectratio=169]{beamer}
\usepackage[english]{babel}
\usepackage[utf8]{inputenc}
\usepackage[T1]{fontenc}
\usepackage{lmodern}
\usepackage{listings}
\usepackage{xcolor}
\usepackage{graphicx}
\usepackage{textcomp}
\DeclareUnicodeCharacter{2014}{---}
\DeclareUnicodeCharacter{2013}{--}
\DeclareUnicodeCharacter{2212}{-}
\DeclareUnicodeCharacter{2192}{\ensuremath{\rightarrow}}
\DeclareUnicodeCharacter{00B1}{\ensuremath{\pm}}
\DeclareUnicodeCharacter{00D7}{\ensuremath{\times}}
\DeclareUnicodeCharacter{00B7}{\ensuremath{\cdot}}
\DeclareUnicodeCharacter{20AC}{EUR}
\DeclareUnicodeCharacter{2070}{\textsuperscript{0}}
\DeclareUnicodeCharacter{00B9}{\textsuperscript{1}}
\DeclareUnicodeCharacter{00B2}{\textsuperscript{2}}
\DeclareUnicodeCharacter{00B3}{\textsuperscript{3}}
\DeclareUnicodeCharacter{2074}{\textsuperscript{4}}
\DeclareUnicodeCharacter{2075}{\textsuperscript{5}}
\DeclareUnicodeCharacter{2076}{\textsuperscript{6}}
\DeclareUnicodeCharacter{2077}{\textsuperscript{7}}
\DeclareUnicodeCharacter{2078}{\textsuperscript{8}}
\DeclareUnicodeCharacter{2079}{\textsuperscript{9}}
\DeclareUnicodeCharacter{2248}{\ensuremath{\approx}}
\DeclareUnicodeCharacter{2264}{\ensuremath{\le}}

% Colors for code
\definecolor{codebg}{RGB}{245,245,245}
\definecolor{codecomment}{RGB}{0,128,0}
\definecolor{codekeyword}{RGB}{0,0,255}
\definecolor{codestring}{RGB}{163,21,21}
\definecolor{bandcolor}{RGB}{200,50,50}

\lstdefinestyle{pythonstyle}{
    language=Python,
    backgroundcolor=\color{codebg},
    basicstyle=\ttfamily\scriptsize,
    keywordstyle=\color{codekeyword},
    commentstyle=\color{codecomment},
    stringstyle=\color{codestring},
    showstringspaces=false,
    breaklines=true,
    frame=single,
    rulecolor=\color{black},
    escapeinside={(*@}{@*)},
    moredelim=**[l][\color{bandcolor}\bfseries]{\#\ ---\ NEW},
    moredelim=**[l][\color{bandcolor}\bfseries]{\#\ ---\ CHANGED},
    literate=
      {á}{{\'a}}1 {é}{{\'e}}1 {í}{{\'i}}1 {ó}{{\'o}}1 {ú}{{\'u}}1
      {Á}{{\'A}}1 {É}{{\'E}}1 {Í}{{\'I}}1 {Ó}{{\'O}}1 {Ú}{{\'U}}1
      {ñ}{{\~n}}1 {Ñ}{{\~N}}1 {ü}{{\"u}}1 {Ü}{{\"U}}1
      {¿}{{?`}}1 {¡}{{!`}}1
      {—}{{---}}1 {–}{{--}}1 {−}{{-}}1
      {±}{{\ensuremath{\pm}}}1 {×}{{\ensuremath{\times}}}1
      {→}{{\ensuremath{\rightarrow}}}1
      {°}{{\textdegree}}1
      {·}{{\ensuremath{\cdot}}}1
      {€}{{\texteuro}}1
      {≈}{{\ensuremath{\approx}}}1
      {≤}{{\ensuremath{\le}}}1
      {⁰}{{\textsuperscript{0}}}1 {¹}{{\textsuperscript{1}}}1
      {²}{{\textsuperscript{2}}}1 {³}{{\textsuperscript{3}}}1
      {⁴}{{\textsuperscript{4}}}1 {⁵}{{\textsuperscript{5}}}1
      {⁶}{{\textsuperscript{6}}}1 {⁷}{{\textsuperscript{7}}}1
      {⁸}{{\textsuperscript{8}}}1 {⁹}{{\textsuperscript{9}}}1
}

\lstset{style=pythonstyle}

% Title slide macro: \lessontitle{Lesson Number}{Title}
\newcommand{\lessontitle}[2]{
    \title{Lesson #1: #2}
    \author{}
    \institute{Tec de Monterrey}
    \begin{frame}[plain]
        \titlepage
    \end{frame}
}

% Figure frame macro: \figureframe{Title}{Image path}
\newcommand{\figureframe}[2]{
    \begin{frame}{#1}
        \begin{center}
            \includegraphics[height=0.8\textheight,width=\textwidth,keepaspectratio]{#2}
        \end{center}
    \end{frame}
}


% Listings pulled from src/ with \lstinputlisting, so the slide is the file.
\lstset{
    basicstyle=\ttfamily\fontsize{8pt}{9.6pt}\selectfont,
    breaklines=true,
    breakatwhitespace=true,
    columns=fullflexible,
    keepspaces=true,
    xleftmargin=0.3em,
    framesep=2pt,
    aboveskip=0pt,
    belowskip=0pt
}

\newcommand{\resultframe}[3]{%
    \begin{frame}{#1}
        \begin{center}
            \includegraphics[height=0.62\textheight,width=\textwidth,keepaspectratio]{#2}
        \end{center}
        \vspace{0.25em}
        {\small #3\par}
    \end{frame}%
}

\newcommand{\claimframe}[2]{%
    \begin{frame}{#1}
        {\small #2\par}
    \end{frame}%
}

\date{}

\begin{document}

\lessontitle{13}{Performance}

\section{This lesson}

\begin{frame}{This lesson}
Lesson 02's evaluation census was a number without a method. This lesson is the method: two counters, one roster, seven steps, same answer.

\vspace{0.6em}
\begin{itemize}
    \item Evaluate once and store: same seed, same \(-87\), 54\% fewer calls.
    \item A cache is worth what the search repeats --- measured, not assumed.
    \item Only a coarser objective actually shrinks the problem, and it must be judged on the real one.
\end{itemize}
\end{frame}

% ================= shifts_01_the_cost_of_a_run =================
\begin{frame}{A work model}
\[
T\approx N_{\mathrm{calls}}W_{\mathrm{call}}+T_{\mathrm{serial}}
+T_{\mathrm{process}}.
\]
A cache is valid only for a deterministic objective keyed by the complete
immutable genome. An exact restart also needs the population, fitness/cache
state, generation counters, and pseudorandom state.
Workers have separate memory under Windows spawn and Unix fork copy-on-write;
worker-local counters and caches are not automatically shared.
\end{frame}

\section{What a run actually costs}

\begin{frame}{What a run actually costs}
This is the last lesson of the course, and it is the only one about the cost of
running a genetic algorithm rather than about the algorithm itself. Everything
that follows is an attempt to make this run cheaper. So the first thing to build
is not an optimisation: it is the meter.
\end{frame}

\resultframe{What a run actually costs}{../figures/shifts_01_the_cost_of_a_run.png}{%
    The run builds \textbf{428} individuals and calls fitness \textbf{930} times — \textbf{2.17 evaluations per individual}, split SELECT 300 / BREED 0 / STATS 630. \textbf{195,300} work units, best fitness \textbf{−87}}

% ================= shifts_02_evaluate_once =================
\section{Calculate the fitness function once}

\begin{frame}{Calculate the fitness function once}
This is the book's section 13.1, and it is the best deal in the whole course.
An individual's genes do not change after it is built, so its fitness cannot
change either. Step 1 nevertheless recomputed it every time a value was needed
- once per individual inside the sort, twice more per individual in the
statistics - and paid a full 210-cell evaluation each time.
\end{frame}

\begin{frame}[fragile]{(1) Individual}
\lstinputlisting[basicstyle=\ttfamily\fontsize{8pt}{9.6pt}\selectfont,firstline=102,lastline=115]{../src/shifts_02_evaluate_once.py}
\end{frame}

\begin{frame}[fragile]{(2) fitness\_of()}
\lstinputlisting[basicstyle=\ttfamily\fontsize{8pt}{9.6pt}\selectfont,firstline=118,lastline=122]{../src/shifts_02_evaluate_once.py}
\end{frame}

\begin{frame}[fragile]{(3) the\_report}
\lstinputlisting[basicstyle=\ttfamily\fontsize{6pt}{7.2pt}\selectfont,firstline=220,lastline=288]{../src/shifts_02_evaluate_once.py}
\end{frame}

\resultframe{Calculate the fitness function once}{../figures/shifts_02_evaluate_once.png}{%
    \textbf{428} calls, \textbf{89,880} work units: \textbf{54.0\%} less, with SELECT and STATS at \textbf{0}. Best fitness still \textbf{−87} — the saving cost nothing}

% ================= shifts_03_the_cache =================
\section{A cache for genomes that come back}

\begin{frame}{A cache for genomes that come back}
Step 2 put the bill at one evaluation per individual built. That is a floor
only for *new* genomes. A genetic algorithm does not produce new genomes all
the time: elites are copied forward, selection picks the same parents twice,
crossover between two copies of one roster returns that roster unchanged, and
a mutation can undo an earlier one. Every one of those is a genome the program
has already priced.

(h=H/(H+M)). If objective work dominates, the ideal work-only speedup is
(1/(1-h)); observed wall-clock speedup also pays lookup and process overhead.
\end{frame}

\begin{frame}[fragile]{(1) CACHE}
\lstinputlisting[basicstyle=\ttfamily\fontsize{8pt}{9.6pt}\selectfont,firstline=97,lastline=103]{../src/shifts_03_the_cache.py}
\end{frame}

\begin{frame}[fragile]{(2) cached\_fitness()}
\lstinputlisting[basicstyle=\ttfamily\fontsize{8pt}{9.6pt}\selectfont,firstline=112,lastline=126]{../src/shifts_03_the_cache.py}
\end{frame}

\begin{frame}[fragile]{(3) Individual}
\lstinputlisting[basicstyle=\ttfamily\fontsize{8pt}{9.6pt}\selectfont,firstline=139,lastline=139]{../src/shifts_03_the_cache.py}
\end{frame}

\begin{frame}[fragile]{(4) the\_report}
\lstinputlisting[basicstyle=\ttfamily\fontsize{6pt}{7.2pt}\selectfont,firstline=244,lastline=310]{../src/shifts_03_the_cache.py}
\end{frame}

\begin{frame}[fragile]{(5) hit\_rate\_scan()}
\lstinputlisting[basicstyle=\ttfamily\fontsize{6pt}{7.2pt}\selectfont,firstline=313,lastline=351]{../src/shifts_03_the_cache.py}
\end{frame}

\resultframe{A cache for genomes that come back}{../figures/shifts_03_the_cache.png}{%
    \textbf{19} hits in 428 builds, \textbf{409} distinct genomes, 4.4\% less work, answer unchanged at \textbf{−87}. The scan shows the hit rate is the *search's* property: \textbf{14.4\%} at mutation p=0.10, \textbf{6.1\%} at 0.25, \textbf{4.4\%} at 0.50, \textbf{3.2\%} at 0.75}

% ================= shifts_04_snapshot =================
\section{Snapshots, and what a snapshot forgets}

\begin{frame}{Snapshots, and what a snapshot forgets}
Long runs get interrupted: a laptop sleeps, a queue job hits its wall time, a
tuning study wants to branch off a run that already cost an hour. The book's
answer (section 13.4) is to write the population's genes to a file and read
them back later.
\end{frame}

\begin{frame}[fragile]{(1) SNAPSHOTS}
\lstinputlisting[basicstyle=\ttfamily\fontsize{8pt}{9.6pt}\selectfont,firstline=60,lastline=64]{../src/shifts_04_snapshot.py}
\end{frame}

\begin{frame}[fragile]{(2) dump\_population()}
\lstinputlisting[basicstyle=\ttfamily\fontsize{8pt}{9.6pt}\selectfont,firstline=152,lastline=158]{../src/shifts_04_snapshot.py}
\end{frame}

\begin{frame}[fragile]{(3) restore\_population()}
\lstinputlisting[basicstyle=\ttfamily\fontsize{8pt}{9.6pt}\selectfont,firstline=161,lastline=166]{../src/shifts_04_snapshot.py}
\end{frame}

\begin{frame}[fragile]{(4) dump\_cache()}
\lstinputlisting[basicstyle=\ttfamily\fontsize{8pt}{9.6pt}\selectfont,firstline=169,lastline=175]{../src/shifts_04_snapshot.py}
\end{frame}

\begin{frame}[fragile]{(5) restore\_cache()}
\lstinputlisting[basicstyle=\ttfamily\fontsize{8pt}{9.6pt}\selectfont,firstline=178,lastline=184]{../src/shifts_04_snapshot.py}
\end{frame}

\begin{frame}[fragile]{(6) the\_report}
\lstinputlisting[basicstyle=\ttfamily\fontsize{6pt}{7.2pt}\selectfont,firstline=283,lastline=382]{../src/shifts_04_snapshot.py}
\end{frame}

\resultframe{Snapshots, and what a snapshot forgets}{../figures/shifts_04_snapshot.png}{%
    30 rosters restore for \textbf{30 evaluations / 6,300 work units} from genes alone, and for \textbf{0 / 0} when the cache is restored first. The cache file is \textbf{14.0×} the size of the gene file — a bad trade at 210 units a call, an excellent one at a minute a call}

% ================= shifts_05_parallel_evaluation =================
\section{Evaluating a generation in a process pool}

\begin{frame}{Evaluating a generation in a process pool}
This is the book's section 13.5, and the first technique in this lesson that
does not remove any work at all. Caching removes evaluations. A pool removes
none: it runs the same evaluations somewhere else, at the same time.
\end{frame}

\begin{frame}[fragile]{(1) Individual}
\lstinputlisting[basicstyle=\ttfamily\fontsize{8pt}{9.6pt}\selectfont,firstline=155,lastline=169]{../src/shifts_05_parallel_evaluation.py}
\end{frame}

\begin{frame}[fragile]{(2) evaluate\_batch()}
\lstinputlisting[basicstyle=\ttfamily\fontsize{8pt}{9.6pt}\selectfont,firstline=172,lastline=182]{../src/shifts_05_parallel_evaluation.py}
\end{frame}

\begin{frame}[fragile]{(3) materialise()}
\lstinputlisting[basicstyle=\ttfamily\fontsize{8pt}{9.6pt}\selectfont,firstline=185,lastline=201]{../src/shifts_05_parallel_evaluation.py}
\end{frame}

\begin{frame}[fragile]{(4) crossover\_operation()}
\lstinputlisting[basicstyle=\ttfamily\fontsize{8pt}{9.6pt}\selectfont,firstline=274,lastline=285]{../src/shifts_05_parallel_evaluation.py}
\end{frame}

\begin{frame}[fragile]{(5) mutation\_operation()}
\lstinputlisting[basicstyle=\ttfamily\fontsize{8pt}{9.6pt}\selectfont,firstline=288,lastline=299]{../src/shifts_05_parallel_evaluation.py}
\end{frame}

\begin{frame}[fragile]{(6) run()}
\lstinputlisting[basicstyle=\ttfamily\fontsize{7pt}{8.5pt}\selectfont,firstline=311,lastline=338]{../src/shifts_05_parallel_evaluation.py}
\end{frame}

\begin{frame}[fragile]{(7) the\_report}
\lstinputlisting[basicstyle=\ttfamily\fontsize{6pt}{7.2pt}\selectfont,firstline=344,lastline=481]{../src/shifts_05_parallel_evaluation.py}
\end{frame}

\resultframe{Evaluating a generation in a process pool}{../figures/shifts_05_parallel_evaluation.png}{%
    Batching alone drops the run from 409 to \textbf{291} evaluations (\textbf{−28.9\%}), because a crossed-then-mutated child used to be evaluated twice. The pooled run gives the identical answer and the identical best-so-far in all 10 generations — while the parent's counters read \textbf{0 evaluations, 0 work units}}

% ================= shifts_06_cache_across_processes =================
\section{Bringing the cache and the meter back across the boundary}

\begin{frame}{Bringing the cache and the meter back across the boundary}
Step 5 finished with a pooled run whose counters read zero evaluations and zero
work units, for a run that had plainly done both. A worker process has separate
memory. Windows spawns it; Unix fork starts with copy-on-write state. In either
case, its cache and counters are not automatically shared with the parent.
\end{frame}

\begin{frame}[fragile]{(1) evaluate\_with\_cost()}
\lstinputlisting[basicstyle=\ttfamily\fontsize{8pt}{9.6pt}\selectfont,firstline=137,lastline=148]{../src/shifts_06_cache_across_processes.py}
\end{frame}

\begin{frame}[fragile]{(2) evaluate\_batch()}
\lstinputlisting[basicstyle=\ttfamily\fontsize{6pt}{7.2pt}\selectfont,firstline=151,lastline=184]{../src/shifts_06_cache_across_processes.py}
\end{frame}

\begin{frame}[fragile]{(3) the\_report}
\lstinputlisting[basicstyle=\ttfamily\fontsize{6pt}{7.2pt}\selectfont,firstline=338,lastline=417]{../src/shifts_06_cache_across_processes.py}
\end{frame}

\resultframe{Bringing the cache and the meter back across the boundary}{../figures/shifts_06_cache_across_processes.png}{%
    Pooled run now reports \textbf{291 / 8 hits / 61,110}, matching the sequential run on all three, answer still \textbf{−87}. Step 5's private per-worker caches had been discarding every hit}

% ================= shifts_07_coarse_fitness =================
\section{A coarse fitness, and what it costs to be cheap}

\begin{frame}{A coarse fitness, and what it costs to be cheap}
Every optimisation so far has been free. Section 13.1 stopped recomputing a
known number; the cache stopped recomputing a known genome; batching stopped
evaluating children that were about to be thrown away; the pool moved work
without changing it. Not one of them altered a single fitness value, and each
one was checked against the previous step's answer to prove it.

For surrogate \(\tilde f\), \(e(x)=\tilde f(x)-f(x)\). Selection may use
\(\tilde f\), but the final candidates must be scored on the real \(f\).
\end{frame}

\begin{frame}[fragile]{(1) coarse\_fitness()}
\lstinputlisting[basicstyle=\ttfamily\fontsize{8pt}{9.6pt}\selectfont,firstline=138,lastline=149]{../src/shifts_07_coarse_fitness.py}
\end{frame}

\begin{frame}[fragile]{(2) evaluate\_with\_cost()}
\lstinputlisting[basicstyle=\ttfamily\fontsize{8pt}{9.6pt}\selectfont,firstline=152,lastline=164]{../src/shifts_07_coarse_fitness.py}
\end{frame}

\begin{frame}[fragile]{(3) the\_report}
\lstinputlisting[basicstyle=\ttfamily\fontsize{6pt}{7.2pt}\selectfont,firstline=352,lastline=448]{../src/shifts_07_coarse_fitness.py}
\end{frame}

\resultframe{A coarse fitness, and what it costs to be cheap}{../figures/shifts_07_coarse_fitness.png}{%
    \textbf{105} work units per call instead of 210, \textbf{51.0\%} less work for the run. The coarse search reaches \textbf{0 on its own objective} — and its roster scores \textbf{−160} on the real one against \textbf{−87}, with \textbf{32} rest violations instead of 17}

\section{Next}

\begin{frame}{Next}
Steps 1--6 all report best fitness \(-87\). The content is the falling bill underneath. Coarse fitness scores 0 on itself and \(-160\) on the real objective.

\vspace{0.8em}
\textbf{The shared Planta F\'isica challenge}

Apply the same discipline there: state the mathematical model, keep feasibility visible, compare under one evaluation currency, and separate observed performance from proof. The current specification lives in the semester repository.
\end{frame}
\end{document}
